% poc/m5-corpus-body.tex — M5 行分割比較用コーパス。各段落をマクロで分離。
% 両エンジンとも fontsize=10pt, zw=10pt, \hsize=200pt（全角20字）で同一条件。
\long\def\PONE{% 区切り約物多め（、。・！？を密に）
近頃の天気は、春らしく、暖かい。しかし、朝方は冷える。桜の便りも、
北から、南から、次々と届く。雨の日は、花びらが、散る。咲く時期は、
土地で違う。早い所では、もう終わり、遅い所では、これからだ。
花見の計画は、天気を見て、決めたい。風の強い日は、散り方がはげしい。}
\long\def\PTWO{% 引用（「」『』）多め
彼は「ほんとうに行くのか」と聞いた。「ああ、行くよ」と私は答えた。
「では、『用意』はできているのか」「うん、『荷物』はまとめた」。
「心配はないのか」「ないわけではないが、『決心』は固い」。
結局、彼は「留守を頼む」と言い、私は「任せておけ」と答えた。}
\long\def\PTHREE{% 括弧類多め
記号の対応は（原則として）次の通りである。括弧は（まず）丸括弧、
〔次に〕角括弧、【強調には】隅付き括弧、《書名には》二重山括弧、
〈短い注には〉山括弧を使う。原文では（引用のため）（併記することも
ある）が、ここでは〔整理のため〕段構造にした。表の列は【見出し】
【量】【単位】の三つで、《付録》を参照されたい。}
